% §10.6 Bibliographic Notes
%
% Monograph register: the four source papers. Fix: d = ceil(ln(1/delta)), natural
% log, matching thm:countmin-bound (the draft wrote "log"). Citations
% cormode-muthukrishnan-2005, flajolet-et-al-2007, broder-1997, indyk-motwani-1998.
% Cross-refs sec:countmin, thm:countmin-bound, sec:hyperloglog,
% prop:hyperloglog-error, sec:minhash-lsh, prop:minhash-error. DEFINED:
% sec:bib-notes-ch10.

\section*{Bibliographic Notes}
\addcontentsline{toc}{section}{Bibliographic Notes}
\label{sec:bib-notes-ch10}

\paragraph{CountMin.}
The construction and the additive bound of \Cref{thm:countmin-bound} are due to
\textcite{cormode-muthukrishnan-2005}: the $d \times w$ table of
pairwise-independent rows, the increment-and-minimum operations, and the
$\hat f_a \le f_a + \epsilon F$ guarantee at $w = \lceil e/\epsilon \rceil$,
$d = \lceil \ln(1/\delta) \rceil$, proved by the Markov-plus-union-bound argument
of \Cref{sec:countmin}. The chapter's addition is the reading of CountMin as a
$\Bernoulli{\mathbb{N}}$ instance, with the upper-triangular channel encoding the
one-sided guarantee and entry-wise merge as composition. Their paper catalogs
applications (heavy hitters, quantiles, inner products) the chapter does not
cover.

\paragraph{HyperLogLog.}
The construction and the $1.04/\sqrt{m}$ relative error of
\Cref{prop:hyperloglog-error} are due to \textcite{flajolet-et-al-2007},
building on the earlier probabilistic-counting work of Flajolet and Martin.
Their generating-function analysis gives the bias correction $\alpha_m$, the
standard-deviation factor, and the limiting log-normal error shape that
underwrites the multiplicative band of the $Q$ matrix; the chapter quotes the
constants without reproducing the derivation. The max-merge and its equivalence
to a sketch of the union are the natural composition under the
max-of-leading-zeros structure.

\paragraph{MinHash and LSH.}
The MinHash sketch is due to \textcite{broder-1997}, whose setting was
near-duplicate web detection; the identity $\Prob{\min h(A) = \min h(B)} = J(A,
B)$ is from that paper, and the $k$-hash estimator and the Hoeffding bound of
\Cref{prop:minhash-error} follow from it. The Binomial-in-$k$ structure of
$\hat J$, which the chapter reads as the $\Bernoulli{[0,1]}$ channel of
\Cref{sec:minhash-lsh}, is already implicit there. Locality-sensitive hashing is
due to \textcite{indyk-motwani-1998}, for approximate nearest-neighbour search;
the MinHash-and-band scheme is the specialization of their general framework to
the Jaccard metric, and their paper also covers Hamming and $\ell_p$ families
for readers wanting the general theory.
